\documentclass[11pt]{article}
\usepackage[a4paper,margin=25mm]{geometry}
\usepackage[T1]{fontenc}\usepackage[utf8]{inputenc}
\usepackage{hyperref}\usepackage{enumitem}
\setlength{\parskip}{0.6em}\setlength{\parindent}{0em}
\begin{document}
\section*{Change letter — CGF fast-track resubmission}
\subsection*{``From Hairballs to Hypotheses: A Survey of AI-Assisted Visual Analytics in the Life Sciences''}
\subsubsection*{EuroVis 2026 STAR submission 1058 $\rightarrow$ Computer Graphics Forum (major revision)}

Authors: Morris Chukhman, Amira Kefi, Nicole M. Chukhman, Silvio Rizzi, Vinayakumar Chalil Karinthayil, Angus Forbes

\bigskip\hrule\bigskip

Dear Editors and Reviewers,

We thank the four reviewers and the EuroVis 2026 STAR committee for the detailed and constructive first-round reviews, and for the fast-track invitation to Computer Graphics Forum. The revision strategy was organized around the Primary Summary's required themes: we rebuilt Section 2 as a fully documented methods section (with a methods supplement now archived as Zenodo v3.1.0), rewrote Sections 3--8 against the fragmentation and balance concerns, added the explicit contribution-positioning section the reviewers asked for (\S3.6), and built the figure program the submission lacked (Figures 2--9; Table 2). Below we address each reviewer concern point by point; where two reviewers raised the same underlying issue, we answer it once and attribute all voices.

\bigskip\hrule\bigskip

\subsection*{1. Motivation and scope: why a new taxonomy? (R1, R3, Primary Summary)}

R1 asked why a new taxonomy is necessary rather than an extension of existing visualization/VA taxonomies; R3 found the necessity of the proposed framework under-justified; the Primary Summary asked for a more precise problem formulation and stronger positioning relative to prior surveys.

\textbf{Response.} The introduction now frames the gap precisely: existing reviews cover representations, tasks, and environments, but ``comparing assistance requires connecting mechanisms and environments to analytic actions'' (\S1). The thesis paragraph states the framework together with the through-line lens — the *locus of analytic agency* — that organizes the entire paper. A new Section 3.6 (``Positioning Against Existing Reviews and Frameworks'') places the work against established agency/automation literature: we cite Endert et al. 2017, Horvitz 1999 as an early mixed-initiative foundation, Brossier et al. 2026 (EuroVis STAR, CGF 45(3)), Ehlers et al., Domova \& Vrotsou, Holter \& El-Assady, Furmanová et al., and Monadjemi et al. The contribution is framed honestly as *applying* established agency and automation concepts to life-science relational analysis — a defensible positioning that R1's and R3's questions, in their different words, were asking for.

\subsection*{2. What counts as ``AI'' in this review (R1, Primary Summary)}

R1 found the boundary of ``AI methods'' unclear — dimensionality reduction has been used for decades; many data-processing methods ``can be interpreted as AI.''

\textbf{Response.} The scope is now stated explicitly and is functional, not chronological: \S1 describes the review as covering ``different modes of computational assistance, rather than successive stages in which newer techniques necessarily provide better assistance''; the \S3 introduction states that ``established clustering, dimensionality reduction, and layout methods qualify by their substantive analytical role, including implementations that do not require learning-based AI.'' The assistance modes are defined by the *role* of computation within interactive analysis, not by technique novelty (\S3.3). Technique-age is deliberately not the inclusion boundary.

\subsection*{3. Methodological transparency (R2, R3, Primary Summary) — the most transformed part of the paper}

R2 asked for a PRISMA-style flowchart with sources and numbers at each stage, a justification of database selection (``surprised that none of the usual databases was used''), and a clearer account of the classification and AI-assisted screening procedures; R3 and the Primary Summary asked for clearer inclusion/exclusion and classification steps.

\textbf{Response.} Section 2 was rebuilt as a full methods section (\S\S2.1--2.7), replacing the submission's \S1.2 and appendix:
\begin{itemize}[leftmargin=*]
\item \textbf{Eligibility} (\S2.1): criteria E1--E5 with stated exclusion boundaries.
\item \textbf{Retrieval} (\S2.2): seven routes — PubMed, Europe PMC, Semantic Scholar, IEEE Xplore, ACM DL, an arXiv snapshot, and prior-survey imports — with the database-coverage question answered in text.
\item \textbf{Deduplication} (\S2.3): 187,446 source occurrences consolidated to 140,959 bibliographic records.
\item \textbf{Screening and candidate classification} (\S2.4, merged for clarity in the final proofing pass): machine-assisted screening outcomes over 140,810 valid records (9,505 Advance / 83,258 Defer / 48,047 Excluded), followed by provisional classification of all 9,505 advanced records; the screening automation is described as provisional coding and its calibration documented.
\item \textbf{Selection and assessment} (\S2.5): six purpose-specific selection sets (20/18/50/200/24/72 = 384) with per-set construction rules; completed full-text assessments now total \textbf{135 systems — 83 eligible, 19 scientifically unresolved, 33 excluded or contextual — with 122 system--cell placements supported across 79 eligible systems} (Figure 3).
\item \textbf{Limitations} (\S2.6): the false-exclusion rate is named as unknown; an independent, stratified human review of the Advance, Defer, and Excluded records is prescribed as the next step; the absence of separate author approval of queue membership before assessment is stated outright.
\item \textbf{Supplementary materials} (\S2.7): the reproducibility package and methods supplement are archived on Zenodo as v3.1.0 (DOI 10.5281/zenodo.23406799); the GenAI-use disclosure appears as a single sentence in \S2.4 (the CGF article class permits no footnotes).
\item \textbf{Figure 1} is the PRISMA-style flow R2 asked for; the methods supplement carries the classification contract, the per-set construction rules, and the new per-system evidence table (\S8 of the supplement: 79 systems, 122 placements, per-row evidence locators).
\end{itemize}

\subsection*{4. Evaluation: paradigms, synthesis, and structured comparison (R1, R2, R3, Primary Summary)}

R1 asked for generalizable evaluation paradigms extracted from the literature; R2 asked for cross-cutting evaluation dimensions across modalities and assistance types; R3 asked for a ``clearer structured comparison of approaches — even if qualitative — to outline strengths and limitations under varying conditions.''

\textbf{Response.} Section 7 was rewritten as evaluation *requirements derived from delegated decisions*: the \S7 introduction's anchor-to-evaluation mapping paragraph derives what each locus of agency must be evaluated on. \textbf{Table 2} is the qualitative structured comparison R3 asked for — nine evaluation examples compared by intervention, study context, evidence, and the limits of the supported claim. Section 7.7 addresses cross-study comparability explicitly, taking the position that heterogeneous studies cannot be pooled but can be compared on what was tested and where the evidence stops (``an unevaluated capability is not evidence of failure'').

\subsection*{5. Visual illustration and the framework overview map (R1, R3, R4, Primary Summary)}

R1 and R3 asked for system images; R4 asked for a comprehensive visual overview mapping surveyed papers onto the framework (``Table 1 only presents a summary of how many papers were found''); the Primary Summary flagged the figure program.

\textbf{Response.} The submission had no system images; the revision has a full figure program:
\begin{itemize}[leftmargin=*]
\item \textbf{Figure 2} — the five interaction environments (with photo credits and an explicit disclaimer that they illustrate interaction environments, not assistance mechanisms).
\item \textbf{Figure 3} — the Assistance $\times$ Visualization Modality map R4 asked for: \textbf{79 systems, 122 supported placements}, with the October access follow-up documented in the caption.
\item \textbf{Figures 4--9} — PeCaX, Facetto ($\times$2), FathomGPT, ProteinSketch, moleculARweb: worked-example images tied directly to the comparisons in \S\S5--6, all reproduced figures carrying CC BY licence statements.
\item The methods supplement's new \textbf{\S8 per-system evidence table} is the tabular companion to Figure 3 (79 systems / 122 placements, per-row evidence locators).
\end{itemize}

\subsection*{6. Structural balance: immersive overemphasis (R2, Primary Summary)}

R2 and the Primary Summary found the immersive environments overemphasized (\S4.3 far more developed than \S\S4.1--4.2; \S5.4 as a special case of adaptive).

\textbf{Response.} The structural fix is the \textbf{limiting rule} (\S3.3): displaying data in VR does not count as immersive assistance; it requires computational use of spatial context. Section 5.4 was rewritten around this rule — placements are now mechanism-level (NAVRNA, ProteinSketch, Vitessce Link's finger-tracking measurement tool) with explicit denials (NUI-VR2, Cosmic View's dome, moleculARweb's display) — making \S5.4 the shortest and most disciplined assistance section. Section 6 was rebalanced across all five environments, and the misplaced \S5.5 (data physicalization/haptics) was removed.

\subsection*{7. Fragmentation and redundancy (R2, Primary Summary)}

R2 flagged ``section fragmentation and redundancies, especially in Sections 5, 6, and 7''; the Primary Summary noted overlap between assistance categories and research challenges.

\textbf{Response.} Sections 4--8 were restructured so each section does one analytic job: tasks (\S4) $\rightarrow$ assistance mechanisms (\S5) $\rightarrow$ modality environments (\S6) $\rightarrow$ evaluation (\S7) $\rightarrow$ challenges (\S8). Section 5.5 was removed; \S8.7 was trimmed to the shared-scenarios infrastructure proposal; the \S2 screening and classification subsections were merged into one.

\subsection*{8. Positioning against prior surveys; missing references (R2-sketch, R4, Primary Summary)}

\textbf{Response — now cited and where:} the immersive-network-analysis distinction is handled in \S3.6; Ceneda et al. 2017 is cited in \S5.2 (definition) and \S8.3, with Ceneda et al. 2019 alongside; Belkacem et al. is cited in \S6.2 per R2's large-display suggestion; Kuťák et al. is cited in \S6.3; VisTrails in \S7.3; Cordeil et al. 2017 (CAVE vs. HMD) in \S6.5, which answers R4 directly; Ens et al.'s immersive-analytics challenges in \S8.6; Shen et al. in \S5.3; Ragan et al. in \S7.3; Preim \& Lawonn in \S\S6.1/7.1; Marriott et al. (eds.) 2018 in \S3.4; Furmanová et al. in \S3.6; Kerren et al. 2017 and O'Donoghue et al. 2018 in \S1; Sperrle et al. and Lam et al. where suggested. Every reference suggested across the reviews is now cited, each placed where suggested. R2's gray-literature sketch items (LLM-networking roadmaps and similar) fall outside the revised framing and were not taken up.

\subsection*{9. Standard terminology: visual analytics vs. visualization (R2, Primary Summary)}

R2 asked that the text use the standard distinction (visual analytics integrates automatic data analysis; otherwise it is visualization). \textbf{Response:} the E3 eligibility criterion and the \S1/\S3 scope language now use the standard distinction throughout.

\subsection*{10. Task-taxonomy references (R2)}

R2 asked that the five task groups be grounded in established task taxonomies. \textbf{Response:} \S3.2 now cites Lee et al. and Brehmer \& Munzner and states explicitly that the five groups are ``developed for this review, not taken directly from either framework.''

\subsection*{11. Internal references; wall-sized displays; adaptive assistance; \S7.5 structure (R2)}

\begin{itemize}[leftmargin=*]
\item \textbf{Internal references} (R2: citation-style [4.1]; R4: use ``see Section 4''; missing whitespace): the draft now uses ``Section X'' prose references throughout, implemented as LaTeX cross-references.
\item \textbf{Wall-sized displays} (R2: \S4.2 shallow on task links; cite Belkacem): \S6.2 was rewritten with task links and coordination design implications; Belkacem et al. is cited.
\item \textbf{Adaptive assistance} (R2: \S5.2 too short, no evaluation considerations, \S5.2/\S5.4 relation): \S5.2 is expanded (Facetto, Metis, SAMIRA; four explicit questions; feedback-loop consequences) and differentiated from \S5.4 by the limiting rule.
\item \textbf{\S7.5 redundancies} (R2: artificial 7.5.1/7.5.2 split): moot — \S\S7--8 were restructured and the old numbering no longer exists.
\end{itemize}

\subsection*{12. Two micro-additions made for this resubmission (R2)}

Two small reviewer-requested references were missing from the revision draft and have been added:
\begin{itemize}[leftmargin=*]
\item \textbf{Semantic-zoom history (R2):} a one-sentence note in the \S4.1 navigation discussion acknowledging the technique's 1990s origins, with R2's two references (10.1145/1456650.1456652; Pook, Lecolinet, Vaysseix \& Barillot, ``Context and interaction in zoomable user interfaces,'' AVI 2000, pp. 227--231, 10.1145/345513.345323).
\item \textbf{Graph backbone extraction (R2):} a general backbone-extraction methods reference (Serrano, Boguñá \& Vespignani, ``Extracting the multiscale backbone of complex weighted networks,'' PNAS 2009) added in \S5.1 alongside the Espadoto et al. dimensionality-reduction survey. The suggested \S5.1.3/\S5.1.1 merge is moot under the restructure.
\end{itemize}

\subsection*{13. Minor corrections (R4 list, R1 minor)}

All verified fixed in the revision: the \S1 comma-before-``that,'' the \S1 double period, Table captions above, and the missing ``??'' reference (the compiled revision contains zero unresolved cross-references).

\bigskip\hrule\bigskip

\subsection*{Closing}

A final proofing pass tightened the front and back of the paper:  the conclusion answers the four research questions explicitly and states how the eligibility criteria bound the claims; \S8's agenda is tied back to RQ2--RQ4.

The core contribution is unchanged from the submission — a survey of AI-assisted visual analytics systems for life-science relational data, organized by an Assistance $\times$ Visualization Modality framework with the locus of analytic agency as its through-line — but the paper that carries it is substantially rebuilt: a documented methods section, an explicit contribution-positioning section, a full figure program, and an evaluation chapter with structured comparisons. We thank the reviewers again for the reviews that made this the paper it is.

On behalf of the authors,

Morris Chukhman, Amira Kefi, Nicole M. Chukhman, Silvio Rizzi, Vinayakumar Chalil Karinthayil, Angus Forbes
15 October 2026

\end{document}
